\subsection{analyse vidéo}

cette sous partie est extrait d'un compte rendu d'un ouvrage collectif sur le projet ASA-SHS, dont voici la présentation :

"Cet ouvrage collectif présente les fondements, la méthodologie et les réalisations techniques d’un vaste projet de recherche mené de 2009 à 2012 à la Fondation Maison des Sciences de l’Homme (FMSH) à Paris par l’Équipe Sémiotique Cognitive et Nouveaux Médias (ESCoM).
Le nom du projet financé par l’ANR (Agence Nationale de la Recherche) Atelier de Sémiotique Audiovisuelle pour les Sciences Humaines et Sociales ; son acronyme est ASA-SHS1. Issu du programme des Archives Audiovisuelles de la Recherche (AAR,
2001- 2009) ce projet a pour ambition de transformer des corpus
audiovisuels scientifiques en ressources de connaissance structurées, interopérables et réutilisables." \footcite{stockinger2011}

Après segmentation, le prolongement est l'analyse de la vidéo et sa description. Stockinger préconise la méthodologie de description suivante : méta description (situe l'analyse et son cadre d'élaboration), la description du document audiovisuel dans son ensemble, et la description des segments spécifiques. \footcite{stockinger2011}

Cette analyse suppose une double exigence : "garantir une standardisation rigoureuse des pratiques analytiques (cohérence documentaire, traçabilité, interopérabilité), tout en respectant la complexité sémiotique du texte audiovisuel, dont la description
doit intégrer plusieurs niveaux de sens, du plus global au plus local." \footcite{stockinger2011}

méta description : 

"positionnement de l’analyste, intentions et contraintes éditoriales, horizon d’usage, ainsi que paramètres institutionnels ou méthodologiques orientant l’acte d’interprétation. [...] elle formalise le point de vue à partir duquel le corpus sera découpé, organisé et décrit". \footcite{stockinger2011}

la méta description permet de décrire les critères de sélection des segments, les objectifs d'analyse, les donénes de gestion \footcite{stockinger2011}

Ensuite, il s'agit d'harmoniser la méta description avec les modèle d'ontologie (ASA-OWL) et que les analyse produites puissent être reliées aux standards du web. \footcite{stockinger2011}

la description du document audiovisuel : 

"il s’agit d’examiner la manière dont
le sens est produit, modulé et orienté par les formes mêmes de la
représentation visuelle et sonore : cadrage, composition, rythme, texture sonore, organisation des voix, relations image/son, etc." \footcite{stockinger2011}

plusiurs plans d'analyse : plan visuel, accoustique, discours. 

le plan visuel : 
analyse de l'image comme un "texte isuel structuré" \footcite{stockinger2011}. 
terme d'isotopie thématique : "la récurrence d’unités visuelles ou de traits sémantiques qui assurent la cohérence d’une séquence ou d’un document." \footcite{stockinger2011}. Cela permet d'identifier les motifs récurrents qui orientent et structurent la lecture de l'image.  
éléments à analyser sur le plan visuel : composition de l'image, mouvements de caméra, échelle des plans, éclairage et couleurs, rythme visuel et montage. 
intéressant : une plateforme de description doit pouvoir expliciter l'identification dun motif technique et son interprétation : aide à relier le sens à la modélisation documentaire. 
lien avec contexte skinsoft : les descriptions extraites peuvent être intégrées à des catégories de champs et de sous champs dans les tables pertinentes. 

analyse du plan accoustique : 
Chion : "la bande-son organise et oriente
la perception du film : elle guide l’attention, structure la temporalité, crée des continuités ou des ruptures, et peut même reconfigurer la lecture de l’image."
le son est un élément structurant de l'image audiovisuelle et crée des effets de sens. 
axes analytiques de la bande son : identification des sources sonores, paramètres techniques, relation son/images, effets expressifs (ambiance, intensité...) \footcite{stockinger2011}

aalyse du discurs : ce que la vidéo énonce, organis,e présuppose comme savoirs : plusieurs éléments : 
les topoï (ensemble de notions, cadres, domaines...)
les sujets de discours 
les thèmes de discours 
\footcite{stockinger2011}

un modèle thématique :
"Un modèle comprend trois couches :
1. Schéma conceptuel : entités, relations, structures topiques.
2. Formulaire d’interface : outil interactif permettant d’opérer la
description.
3. Instances de description : analyses produites par les utilisateurs." \footcite{stockinger2011}

"chaque nouveau modèle peut être
enrichi, adapté ou dérivé en fonction des besoins spécifiques d’un corpus." \footcite{stockinger2011}

dernière partie de l'analyse d'une vidéo : usage :
"la manière dont une ressource audiovisuelle peut être utilisée, adaptée et diffusée dans des contextes variés" \footcite{stockinger2011}
éducatif, patrimonial, politique, social...

"Une vidéo non segmentée possède une identité auctoriale (c’est-à-dire liée à son auteur, individuel ou collectif), mais le découpage effectué dans l’Atelier de Segmentation (chapitre 2) permet de multiplier les usages possibles : chaque segment devient un support susceptible d’être exploité différemment selon le public visé." \footcite{stockinger2011}

la segmentation a plutot une visée pédagogique : isoler des extraits exploitables. 

Traduction multilingue : pouvoir ajouter des versions traduites des segments. 

la description thématiques peut permettre à terme la recherche par thématiques qui donnerait tous les segments reliés à cette thématique sur une interface web. 

combiner lecture audiovisuelle et exploration documentaire : les résultats de recherche affichent la vidéo, sa description, ses éléments associés, et suggestion de vidéos pertinentes (même thème, même auteur, époque...)

Accès thématique par thésaurus (quel thésaurus cinéma utilisé par skinsoft ?) voir appli albert kahn 

importance de la contextualisation d'une vidéo en tout temps : doit êgtre rattachée à un thème, un projet, une description, un évènement... 

OntoEditor : outil d'édition d'ontologies + formalise les connaissances : strcture les savoirs de l'analyse audiovisuelles + garantit leur cohérence, interopréabilité et réutilisation

application INterview : structure les descriptions qui articulent les observations analytiues sur les segments par ex. 

applications Internet riches (RIA) : notion intéressantes qui priorise une navigtaion fluide et interactive 

produire une carte conceptuelle du processus d'analyse et de description. 